% 02-introduction.tex
\chapter{Einleitung}
\label{chap:introduction}

Dieses Kapitel dient als Einführung in die Thematik der wissenschaftlichen Arbeit. Es soll den Lesenden abholen, die Relevanz des Themas verdeutlichen und die Forschungsfragen sowie den Aufbau der Arbeit darlegen.

Als Beispiel für die Verwendung von Abkürzungen: Die Forschung im Bereich der \ac{AI} hat in den letzten Jahren stark zugenommen. Bei der ersten Nennung wird die Abkürzung automatisch ausgeschrieben, danach nur noch als \ac{AI} bezeichnet. Ein weiteres Beispiel für den Plural sind \acp{svm}.

\section{Hintergrund und Motivation}
\label{sec:motivation}
Hier wird der allgemeine Hintergrund erläutert. Warum ist das Thema relevant? 
Ein grundlegendes Verständnis der Softwareentwicklung ist essenziell \cite{meyer2020grundlagen}. Durch Zitationen belegst du deine Aussagen und verweist auf die verwendete Literatur. Manchmal ist es auch nützlich, zusätzliche Erklärungen in einer Fußnote auszulagern\footnote{Dies ist ein Beispiel für eine Fußnote am unteren Rand der Seite.}.

\section{Zielsetzung und Forschungsfragen}
Was soll mit dieser Arbeit erreicht werden? Welche Fragen sollen beantwortet werden, basierend auf der Motivation aus Abschnitt~\ref{sec:motivation}?
\begin{itemize}
  \item \textbf{Ziel 1:} Entwicklung einer neuen Methode zur Bilderkennung.
  \item \textbf{Ziel 2:} Evaluation der Methode anhand eines bestehenden Datensatzes.
  \item \textbf{Ziel 3:} Bereitstellung einer sauberen Dokumentation.
\end{itemize}

\section{Aufbau der Arbeit}
Die Arbeit gliedert sich wie folgt:
Kapitel~\ref{chap:related_work} beleuchtet den aktuellen Stand der Technik. In Kapitel~\ref{chap:methodology} wird die Methodik vorgestellt, gefolgt von der praktischen Implementierung in Kapitel~\ref{chap:implementation}. Abschließend fasst Kapitel~\ref{chap:conclusion} die Arbeit zusammen und gibt einen Ausblick auf zukünftige Forschungsmöglichkeiten.